\chapter{Electrical Conductivity: What DFT Can and Cannot Tell You}
\label{ch:conductivity}
% generated by verification/optics_report.py -- do not edit
\newcommand{\GF}{0.264}
\newcommand{\GFb}{0.276}
\newcommand{\VF}{1.33}
\newcommand{\VFb}{1.31}
\newcommand{\WP}{8.47}
\newcommand{\WPb}{8.55}
\newcommand{\ST}{1.47}
\newcommand{\STb}{1.49}
\newcommand{\NEFF}{0.90}
\newcommand{\TAU}{30}
\newcommand{\MFP}{40}
\newcommand{\LAMTR}{0.14}
\newcommand{\NKFINE}{1661}
\newcommand{\NKCOARSE}{413}
\newcommand{\NTRANS}{199}
\newcommand{\EDGE}{1.9}
\newcommand{\ONSET}{1.5}
\newcommand{\EDGES}{2.1}
\newcommand{\EPSZERO}{9.9}
\newcommand{\NEFFINTER}{1.4}
\newcommand{\RGBLDA}{195,182,177}
\newcommand{\RLDAb}{0.44}
\newcommand{\RLDAg}{0.46}
\newcommand{\RLDAr}{0.65}
\newcommand{\RGBSHa}{243,162,157}
\newcommand{\RSHab}{0.36}
\newcommand{\RSHag}{0.38}
\newcommand{\RSHar}{0.84}
\newcommand{\RGBSHb}{255,180,129}
\newcommand{\RSHbb}{0.26}
\newcommand{\RSHbg}{0.52}
\newcommand{\RSHbr}{0.92}
\newcommand{\RGBDR}{253,253,253}
\newcommand{\RDRb}{0.98}
\newcommand{\RDRg}{0.98}
\newcommand{\RDRr}{0.98}
\newcommand{\RGBCO}{179,184,168}
\newcommand{\RCOb}{0.36}
\newcommand{\RCOg}{0.44}
\newcommand{\RCOr}{0.63}
\newcommand{\RGBETA}{205,183,175}
\newcommand{\RETAb}{0.44}
\newcommand{\RETAg}{0.48}
\newcommand{\RETAr}{0.69}
\definecolor{goldLDA}{RGB}{195,182,177}
\definecolor{goldSHa}{RGB}{243,162,157}
\definecolor{goldSHb}{RGB}{255,180,129}
\definecolor{goldDR}{RGB}{253,253,253}
\newcommand{\RGBEXP}{255,222,163}
\definecolor{goldEXP}{RGB}{255,222,163}
\newcommand{\REXPb}{0.41}
\newcommand{\REXPg}{0.79}
\newcommand{\REXPr}{0.96}
\newcommand{\EDGEEXP}{2.4}
\newcommand{\EPSEXP}{5.7}
\newcommand{\EPSCALC}{8.3}
\newcommand{\EDGENOSOC}{2.04}
\newcommand{\EDGESOC}{1.85}


Gold is one of the best conductors there is: at room temperature its
resistivity is $\rho = 2.2\,\mu\Omega\,$cm ($22$\,n$\Omega$\,m at
$293$\,K), i.e.\ a conductivity
$\sigma = 1/\rho = 4.5\times10^{7}$\,S/m. Can the program you have written
predict this number? The honest answer has two halves, and separating them
is the point of this chapter.
\begin{itemize}[leftmargin=2em]
\item \textbf{Yes:} everything about conduction that is a property of the
\emph{perfect crystal} --- which electrons carry the current, how fast they
move, how strongly they respond to a field --- follows from the bands and
wavefunctions you already have. These combine into one number, the
\emph{plasma frequency} $\omega_p$, and you will compute it.
\item \textbf{Not by itself:} in a perfect, rigid crystal the conductivity
is \emph{infinite}. The measured value is finite only because electrons are
scattered --- at room temperature by the thermal vibrations of the lattice
--- and the scattering rate is not a ground-state property of the frozen
crystal. It \emph{can} be computed from first principles, but that needs
the lattice vibrations and their coupling to the electrons, a substantial
extension of the code that is outlined, not built, in
Section~\ref{sec:tau}.
\end{itemize}
The one new tool introduced here, the velocity operator, is also the key to
the optical properties of the next chapter.

\section{Why a perfect crystal does not resist}

A Bloch state $\psi_{n\kk}$ is a stationary state of the periodic
potential: an electron in it moves through the lattice forever with the
mean velocity (derived in the next section)
\begin{equation}
\mathbf v_{n\kk} = \pdv{\varepsilon_{n\kk}}{\kk}
\label{eq:groupvel}
\end{equation}
--- the group velocity of a wave packet built from states near $\kk$. The
periodic potential of the ions, however strong, does not scatter it; that
is the content of Bloch's theorem. In equilibrium the currents of all
occupied states cancel, because $\varepsilon_{n,-\kk} = \varepsilon_{n\kk}$
implies $\mathbf v_{n,-\kk} = -\mathbf v_{n\kk}$.

Now apply a uniform electric field $\mathbf E$. The semiclassical equation
of motion of a Bloch electron (charge $-1$ in atomic units) is
\begin{equation}
\frac{\dd\kk}{\dd t} = -\mathbf E :
\end{equation}
the force changes the \emph{crystal momentum}, so every electron drifts
through $\kk$-space at the same rate. The occupied region of the Brillouin
zone shifts as a whole, the cancellation between $\kk$ and $-\kk$ is lost,
and a current flows that keeps growing as long as the field is on. Nothing
in the perfect crystal stops it. A finite, steady current requires
processes that knock electrons back toward equilibrium: collisions with
lattice vibrations (phonons), with impurities and defects, and with each
other. Summarize them, crudely but usefully, by a single \emph{relaxation
time} $\tau$: the average time an electron travels between collisions.

\section{The velocity operator}
\label{sec:velocity}

Both this chapter and the next need matrix elements of the electron
velocity between Bloch states. At fixed $\kk$ your program represents the
Hamiltonian as a matrix $H_{\kk}$ acting on the coefficients $c(\GG)$ ---
the Hamiltonian for the periodic functions $u_{n\kk}$. The velocity
operator in the same representation is simply its derivative,
\begin{equation}
\boxed{\;\hat{\mathbf v} = \pdv{H_{\kk}}{\kk}\;}
\qquad\qquad
\mathbf v_{mn}(\kk) = \sum_{\GG\GG' ss'} c^{s\,*}_{m\kk}(\GG)\,
\pdv{H^{ss'}_{\GG\GG'}(\kk)}{\kk}\, c^{s'}_{n\kk}(\GG') .
\label{eq:velop}
\end{equation}
\emph{Why.} Differentiate the eigenvalue $\varepsilon_{n\kk} =
\braket{u_{n\kk}|H_\kk|u_{n\kk}}$ with respect to $\kk$. Three terms
appear; the two containing $\partial u/\partial\kk$ add up to
$\varepsilon_{n\kk}\,\partial_\kk\braket{u_{n\kk}|u_{n\kk}} = 0$ because
the state stays normalized. What remains is the
\textbf{Hellmann--Feynman theorem},
\begin{equation}
\pdv{\varepsilon_{n\kk}}{\kk} = \Braket{u_{n\kk}|\pdv{H_\kk}{\kk}|u_{n\kk}}
= \mathbf v_{nn}(\kk),
\label{eq:hf}
\end{equation}
which is Eq.~\eqref{eq:groupvel}: the diagonal matrix element of
$\partial H_\kk/\partial\kk$ is the group velocity. The off-diagonal
elements $\mathbf v_{mn}$, $m\ne n$, govern transitions between bands
(Chapter~\ref{ch:colour}).

\paragraph{Term by term.} Go through the Hamiltonian of
Chapter~\ref{ch:hamiltonian} and ask where $\kk$ appears.
\begin{itemize}[leftmargin=2em]
\item \emph{Kinetic energy} $\tfrac12|\kk+\GG|^2$: derivative $\kk+\GG$.
This is the momentum, as for a free particle.
\item \emph{Local potential} $V_{\mathrm{eff,loc}}(\GG-\GG')$: no $\kk$
dependence, no contribution.
\item \emph{Nonlocal pseudopotential}: the projectors
$\tilde\beta_p(\kk+\GG)$ of Eq.~\eqref{eq:betaG} depend on $\kk$, so this
term \emph{does} contribute. (Physically: a nonlocal potential does not
commute with the position operator, so velocity is no longer momentum
alone. Leaving this term out is the classic error in pseudopotential
optics codes.)
\end{itemize}
Writing $\mathbf q = \kk+\GG$ and $\partial_\alpha\tilde\beta_p \equiv
\partial\tilde\beta_p(\mathbf q)/\partial q_\alpha$ for the Cartesian
directions $\alpha = x,y,z$, and with the projections
$P_{n} = \sum_\GG\tilde\beta^*(\mathbf q)\,c_n(\GG)$ of
Algorithm~\ref{alg:applyh} and their analogues
$\partial_\alpha P_{n} = \sum_\GG[\partial_\alpha\tilde\beta(\mathbf q)]^*
c_n(\GG)$ (both are vectors over the composite index: projector $\times$
spin),
\begin{equation}
v^{\alpha}_{mn} = \sum_{\GG s} c^{s\,*}_{m}(\GG)\,(\kk+\GG)_\alpha\,
c^{s}_{n}(\GG)
\;+\; (\partial_\alpha P_m)^\dagger D\, P_n
\;+\; P_m^\dagger D\, (\partial_\alpha P_n).
\label{eq:vmn}
\end{equation}
The derivative of a projector is the derivative of an explicit function,
$\tilde\beta_p(\mathbf q) \propto \tilde p_{i\ell}(|\mathbf q|)\,
Y_{\ell m}(\hat{\mathbf q})$: obtain it with automatic differentiation
(\texttt{jax.jacfwd} of the function that returns the projector table for
given $\mathbf q$), or with a central difference
$[\tilde\beta(\mathbf q + \delta\,\hat{\mathbf e}_\alpha)
- \tilde\beta(\mathbf q - \delta\,\hat{\mathbf e}_\alpha)]/2\delta$,
$\delta = 10^{-4}$, which is accurate to $10^{-8}$ for these smooth
functions.

\begin{algorithm}[H]
\caption{Velocity matrix elements at one $\kk$-point}
\label{alg:velocity}
\begin{algorithmic}[1]
\Function{Velocity}{$C[n,s,\GG]$;\ basis}
\State $P[n,(s,p)] \gets \sum_\GG \tilde\beta_p^*(\mathbf q)\,C[n,s,\GG]$;\quad
$DP \gets P\,D^{\mathsf T}$ \Comment{as in \textsc{ApplyH}}
\For{$\alpha = x, y, z$}
\State $\partial_\alpha\tilde\beta_p(\mathbf q) \gets$ derivative of the
projector table with respect to $q_\alpha$
\State $\partial P[n,(s,p)] \gets \sum_\GG
[\partial_\alpha\tilde\beta_p(\mathbf q)]^*\,C[n,s,\GG]$
\State $v^\alpha_{mn} \gets \sum_{\GG s}\overline{C[m,s,\GG]}\; q_\alpha\, C[n,s,\GG]$
\Comment{kinetic part}
\State $v^\alpha_{mn} \mathrel{+}= \sum_{a}\overline{\partial P[m,a]}\,DP[n,a]
+ \sum_{a}\overline{P[m,a]}\,(\partial P\,D^{\mathsf T})[n,a]$
\Comment{nonlocal part}
\EndFor
\State \Return $v^\alpha_{mn}$ \Comment{three Hermitian $N_b\times N_b$ matrices}
\EndFunction
\end{algorithmic}
\end{algorithm}

\paragraph{Degenerate bands.} Every band of gold is a Kramers pair, and
within a degenerate pair the eigensolver returns an arbitrary basis. The
group velocity of the pair is therefore not an individual diagonal element
but the $2\times2$ block of $v^\alpha$ within the pair. For gold the
combined symmetry $\hat P\Theta$ (Section~\ref{sec:symmetry}) forces that
block to be a multiple of the unit matrix: both partners move with the
same velocity, and the block is the same in any basis. Use
$\tfrac12\mathrm{tr}$ of the block for the velocity and
$\tfrac12\sum_{m,n\in\text{pair}}|v^\alpha_{mn}|^2$ for its square.

\paragraph{Test.} Equation~\eqref{eq:hf} is a complete check of
Algorithm~\ref{alg:velocity}, nonlocal term included: solve at
$\kk \pm h\,\hat{\mathbf e}_\alpha$ ($h = 10^{-4}$\,Bohr$^{-1}$) and compare
$[\varepsilon_n(\kk+h\hat{\mathbf e}_\alpha) -
\varepsilon_n(\kk-h\hat{\mathbf e}_\alpha)]/2h$ with the pair-averaged
$v^\alpha_{nn}$. In the reference calculation they agree to better
than $10^{-6}$\,a.u.\ for all bands at a generic $\kk$, the matrices are
Hermitian to rounding error, and the off-diagonal element inside each
Kramers pair is below $10^{-8}$. Without the nonlocal term the
discrepancy is $0.1$--$0.3$\,a.u.\ for every band --- as large as the
velocities themselves.

\section{Conductivity in the relaxation-time approximation}

In a steady field, an electron accumulates the shift $\delta\kk = -\mathbf
E\,\tau$ before it is scattered. The occupation of state $n\kk$ is then
that of the equilibrium distribution $f^0$ at the point the electron
started from,
\begin{equation}
f_{n\kk} = f^0(\varepsilon_{n,\kk+\mathbf E\tau})
\approx f^0(\varepsilon_{n\kk})
+ \pdv{f^0}{\varepsilon}\;\mathbf v_{n\kk}\cdot\mathbf E\,\tau .
\end{equation}
The current density is the sum of charge ($-1$) times velocity over all
occupied states, per unit volume; with our conventions (one electron per
spinor state, weights $\sum_\kk w_\kk = 1$, cell volume $\Omega$) that is
$\mathbf j = -\frac{1}{\Omega}\sum_{n\kk}w_\kk f_{n\kk}\mathbf v_{n\kk}$.
The equilibrium part cancels, and what is left is linear in the field,
$j_\alpha = \sum_\beta\sigma_{\alpha\beta}E_\beta$, with
\begin{equation}
\sigma_{\alpha\beta} = \frac{\tau}{\Omega}\sum_{n\kk} w_\kk
\Big(-\pdv{f^0}{\varepsilon}\Big)_{\varepsilon_{n\kk}}
v^\alpha_{n\kk}\,v^\beta_{n\kk} .
\label{eq:sigmatensor}
\end{equation}
The factor $-\partial f^0/\partial\varepsilon$ is a narrow peak at the
Fermi level (a $\delta$-function at zero temperature): \emph{only
electrons at the Fermi surface conduct}. In a cubic crystal the tensor is
isotropic, $\sigma_{\alpha\beta} = \sigma\,\delta_{\alpha\beta}$, and it is
customary to write
\begin{equation}
\boxed{\;
\sigma = \frac{\omega_p^2\,\tau}{4\pi},
\qquad
\omega_p^2 = \frac{4\pi}{\Omega}\sum_{n\kk} w_\kk\,
\delta(\varepsilon_{n\kk}-\mu)\;\frac{|\mathbf v_{n\kk}|^2}{3}\;}
\label{eq:plasma}
\end{equation}
(atomic units). $\omega_p$ is the \textbf{Drude plasma frequency}. It
packages everything the band structure contributes to conduction, and it
is a ground-state quantity your code can deliver.

\paragraph{Sanity check: free electrons.} For a free-electron gas of density
$n$, with $g(\varepsilon_F) = 3n/2\varepsilon_F$ states per unit energy
and volume and $v_F^2 = 2\varepsilon_F$, Eq.~\eqref{eq:plasma} gives
$\omega_p^2 = \tfrac{4\pi}{3}\,g(\varepsilon_F)\,v_F^2 = 4\pi n$ and
$\sigma = n\tau$ --- in SI units the Drude formula $\sigma = ne^2\tau/m$
of elementary solid-state physics. For gold with one $6s$ electron per
atom, $n = 1/\Omega$: $\hbar\omega_p = 9.02$\,eV and
$v_F = (3\pi^2 n)^{1/3} = 1.39\times10^6$\,m/s. The $d$ electrons, whose
bands lie entirely below $E_F$, do not conduct.

\paragraph{Evaluating the Fermi-surface sum.} The $\delta$-function is
replaced by a normalized Gaussian of width $s$,
$\delta_s(x) = \ee^{-x^2/s^2}/(s\sqrt\pi)$, and the sum is taken over a
mesh much denser than the SCF mesh, using \textsc{SolveBands} in the fixed
self-consistent potential (Section~\ref{sec:nscf}). The integrand is
invariant under the 48 operations of the cubic point group (because it
contains $|\mathbf v|^2$, not individual components), so it suffices to
visit the irreducible wedge: group the mesh points into sets related by
the 48 signed permutations of the Cartesian components, keep one point of
each set, and give it the weight (size of the set)$/N_\kk$. For a
$24^3$ mesh that is 413 points instead of 13\,824.

\begin{algorithm}[H]
\caption{Drude weight, density of states and Fermi velocity}
\label{alg:drude}
\begin{algorithmic}[1]
\State irreducible $\kk$-points and weights $w_\kk$ of an $N^3$
$\Gamma$-centered mesh
\For{each $\kk$}
\State $(\varepsilon, C) \gets$ \Call{SolveBands}{$\kk$};\quad
$v^\alpha_{mn} \gets$ \Call{Velocity}{$C$}
\State $V_n \gets \tfrac13\cdot\tfrac12\sum_\alpha\sum_{m,m'\in
\text{pair}(n)}|v^\alpha_{mm'}|^2$ \Comment{$|\mathbf v_{n\kk}|^2/3$}
\State store $\varepsilon_{n\kk}$, $V_{n\kk}$, and
$\tfrac13\sum_\alpha|v^\alpha_{mn}|^2$ for all pairs $m,n$
\Comment{the last for Chapter~\ref{ch:colour}}
\EndFor
\State $g(\mu) \gets \sum_{n\kk}w_\kk\,\delta_s(\varepsilon_{n\kk}-\mu)$
\Comment{states per Ha per cell}
\State $\omega_p^2 \gets \dfrac{4\pi}{\Omega}\sum_{n\kk}w_\kk\,
\delta_s(\varepsilon_{n\kk}-\mu)\,V_{n\kk}$;\quad
$\langle v_F^2\rangle \gets 3\sum w\,\delta_s V\big/\sum w\,\delta_s$
\end{algorithmic}
\end{algorithm}

\section{Results for gold}

Reference calculation: fixed potential of Case~B (Chapter~\ref{ch:verify};
LDA, spin--orbit coupling on), $\Gamma$-centered meshes reduced to the
irreducible wedge.
\begin{center}
\begin{tabular}{@{}lcccc@{}}
\toprule
 & \multicolumn{2}{c}{this calculation} & free electrons & experiment\\
 & $24^3$ & $40^3$ & ($6s^1$) & \\
\midrule
$g(E_F)$ (states/eV/atom) & \GF & \GFb & $0.27$ & $\approx 0.3$\\
r.m.s.\ Fermi velocity ($10^6$\,m/s) & \VF & \VFb & $1.39$ & $\approx 1.4$\\
$\hbar\omega_p$ (eV) & \WP & \WPb & $9.02$ & $8.45$--$9.0$\\
$\sigma/\tau$ ($10^{6}$\,S\,m$^{-1}$\,fs$^{-1}$) & \ST & \STb & $1.66$ & ---\\
\bottomrule
\end{tabular}
\end{center}
(Gaussian width $s = 0.005$\,Ha. On the $40^3$ mesh, $s = 0.01$\,Ha
raises $\hbar\omega_p$ by $0.2$\,eV and $s = 0.003$\,Ha lowers it by
$0.5$\,eV: so narrow a Gaussian samples the Fermi surface unevenly on a
mesh of this density. Treat the numbers as accurate to a few tenths of an
eV. Experimental values:
$g(E_F)$ from the electronic specific heat coefficient
$\gamma = 0.69$\,mJ\,mol$^{-1}$K$^{-2}$, which includes an enhancement
by the electron--phonon interaction of the order of 10\%; $\omega_p$
from Drude fits to optical data, $8.45$\,eV (Olmon \emph{et al.}, 2012)
to $9.0$\,eV (Johnson and Christy, 1972).) Three
things to take from the table. First, gold conducts like a free-electron
metal with close to one electron per atom: the effective number of
carriers, $\omega_p^2\Omega/4\pi$, comes out as \NEFF{}. The single band
crossing $E_F$ in Fig.~\ref{fig:bands} is the $6s$ band. Second,
Kohn--Sham theory is on firm ground here: unlike the position of the $d$
bands, the Fermi surface of a simple metal is described well. Third, the
band structure alone fixes the ratio $\sigma/\tau$ --- but not $\sigma$.

\paragraph{Unit conversions.} One atomic unit of conductivity is
$e^2/(a_0\hbar) = 4.5998\times10^{6}$\,S/m; of time,
$\hbar/\mathrm{Ha} = 0.024189$\,fs; of velocity,
$2.1877\times10^{6}$\,m/s. Example: $\hbar\omega_p = 8.5$\,eV is
$\omega_p = 0.3124$\,a.u.; with $\tau = 30$\,fs $= 1240$\,a.u.,
$\sigma = \omega_p^2\tau/4\pi = 9.63$\,a.u.\ $= 4.4\times10^7$\,S/m.

\section{The missing ingredient: the relaxation time}
\label{sec:tau}

\paragraph{From experiment.} Turn the argument around: the measured
room-temperature resistivity and the computed plasma frequency give
\begin{equation}
\tau = \frac{4\pi\sigma}{\omega_p^2} \approx \TAU\,\mathrm{fs},
\qquad
\ell = v_F\,\tau \approx \MFP\,\mathrm{nm} .
\end{equation}
An electron at the Fermi surface of gold travels about a hundred lattice
constants between collisions: the picture of nearly free Bloch waves,
occasionally scattered, is self-consistent. (This is also why thin gold
films and nanowires conduct worse than bulk gold: once the sample is not
much thicker than $\ell$, surface scattering matters.)

\paragraph{From first principles (outline only).} At room temperature
$\tau$ in a pure metal is set by electron--phonon scattering. Above the
Debye temperature (about 165\,K for gold) the result of the full theory is
remarkably simple,
\begin{equation}
\frac{1}{\tau} = 2\pi\,\lambda_{\mathrm{tr}}\,k_BT
\qquad\Longrightarrow\qquad
\rho(T) = \frac{4\pi}{\omega_p^2}\cdot 2\pi\,\lambda_{\mathrm{tr}}\,k_BT ,
\label{eq:rhoT}
\end{equation}
(atomic units; $k_BT = 9.34\times10^{-4}$\,Ha at 295\,K): the resistivity
is linear in temperature, as observed, with a slope fixed by $\omega_p$
and one dimensionless number, the \emph{transport electron--phonon
coupling constant} $\lambda_{\mathrm{tr}}$. Inserting the measured
resistivity and the computed $\omega_p$ gives
$\lambda_{\mathrm{tr}}\approx\LAMTR$ --- a small value: gold is a good
conductor because its electrons couple weakly to its lattice vibrations
(the same reason it is not a superconductor at any accessible
temperature). This is exactly how ``empirical'' coupling constants are
extracted in the literature: Savrasov and Savrasov (1996, Eq.~23 there)
use $\lambda_{\mathrm{tr}} = c_1\omega_p^2/8\pi^2k_B$ with $c_1$ the
measured slope of $\rho(T)$, which is Eq.~\eqref{eq:rhoT} solved for
$\lambda_{\mathrm{tr}}$, and find $0.12$ for copper (and values above 1
for good superconductors such as lead and niobium).

Computing $\lambda_{\mathrm{tr}}$ is a research-grade but well-defined DFT
calculation. It requires, beyond this document: (i) forces on atoms
(derivatives of the total energy with respect to atomic positions --- in
JAX, \texttt{grad} of your energy function); (ii) the phonon frequencies
and displacement patterns, from the forces in supercells with one atom
displaced, or from density-functional perturbation theory; (iii) the
change of the self-consistent potential caused by each phonon, and its
matrix elements between Bloch states at the Fermi surface; (iv) a double
Fermi-surface integral of those matrix elements weighted by the change of
velocity in each scattering event. Every step uses machinery you now
understand --- supercells, fixed-potential band calculations, velocity
matrix elements, Fermi-surface sums --- but none of it has been carried
out for this document, and Eq.~\eqref{eq:rhoT} is quoted from the
literature, not derived. That the programme works is established: a
published first-principles Boltzmann calculation for gold (Gunst, Blom
and Stokbro, 2018) obtains $27.0$\,n$\Omega$\,m at 300\,K, some $20$--$30\%$ above the
measured room-temperature value ($\approx 22$\,n$\Omega$\,m), and reproduces the ordering Au $>$ Cu $>$ Ag. The residual resistivity at low temperature is different again:
it is caused by impurities and defects and is a property of the sample,
not of gold.

\begin{keybox}{Summary}
The DC conductivity factorizes, $\sigma = \omega_p^2\tau/4\pi$. The plasma
frequency $\omega_p$ is a Fermi-surface average of squared group
velocities, computed from $\hat{\mathbf v} = \partial H_\kk/\partial\kk$
(kinetic \emph{and} nonlocal-pseudopotential terms) in a fixed-potential
calculation on a dense mesh; for gold it is that of one nearly free
electron per atom. The relaxation time $\tau$ is not a property of the
perfect frozen crystal: take it from experiment, or compute the
electron--phonon coupling --- a separate project.
\end{keybox}

\begin{gpubox}{GPU/JAX notes}
The $\kk$-points of the dense mesh are independent: this is
\textsc{SolveBands} batched over $\kk$ (\texttt{vmap}, in chunks), followed
by three \texttt{einsum}s per $\kk$ for the velocity matrices. The
projector derivative is one \texttt{jax.jacfwd} of a function you already
have; because it is evaluated once per $\kk$-point, its cost is
irrelevant. Store $\varepsilon_{n\kk}$ and the real arrays
$\tfrac13\sum_\alpha|v^\alpha_{mn}|^2$ --- a few megabytes --- and do all
of the analysis of this chapter and the next on the host.
\end{gpubox}
